% v2 chapter. The method text of the paper is input; thesis-only parts are the design rationale, the exploratory screening (numbers copied from
% DECISIONS.md, all development data) and the thesis-length reading of the dose-response, ablation and encoder results (generated tables).
\section{Why a second analysis plan}
The first study left three things open (Chapter~\ref{ch:v1}). Its models kept lesions on rendered data but not on real radiographs; lesion
insertion acted only through the supervised target, so the network was never told that the bone estimate should not change when a lesion is
added; and the comparison with the reference model on SZCH-X-Rays had been looked at, so that set could no longer decide anything about a new
recipe. The second plan, SAP\_v2.md, was written on 29 October 2026 after the new recipe had been developed on development data only: the 24
SZCH development pairs, the audit on all 75 SZCH images, and the synthetic validation data. It names five confirmatory hypotheses on sets that no
decision had used: the dose-response of absorption on the twin weight on VinDr-RibCXR (H1$'$), non-inferior fidelity on the 51 held-out SZCH
pairs (H2$'$), lower absorption than the reference model and the no-twin control on VinDr-RibCXR (H3$'$), evidence that does not use insertion
(H4$'$: real-nodule contrast and downstream tasks), and reliability (H5$'$). The plan, the configuration and the weights were frozen on
1~November 2026 (\path{freeze_v2.json}) before any of these sets was evaluated.

\section{Training on public CT projections with a twin regulariser}
\label{sec:method_v2}
% PLAN (about 1.5 pages). This section is new text for the v2 recipe; methods.tex (v1 recipe, physics, audit definition) is untouched.
%   4.1 Qure projections as an exactly additive source: what a row contains, why the bone image is exact (composite = a*soft + b*bone + c*lung + d),
%       the augmentation order of s4/qp.py, and why monotone tone curves keep the decomposition exact.
%   4.2 Twin regulariser: definition, loss, the weight w, and the variants that form the ablation (stop-gradient, support size, whole image, number
%       of pasted lesions, insertion without the tie term). Relation to the v1 counterfactual insertion (CLIA) in methods.tex.
%   4.3 Ensemble and uncertainty: K=3 members, mean bone image, spread, split-conformal scale, what is calibrated where (SAP_v2.md, H5').
%   4.4 Training budget and what is identical across arms (same optimiser and schedule as methods.tex).
% Verify every formula below against s4/train.py, s4/qp.py and s4/uncertainty.py at the freeze commit before submission.
\subsection{Qure projections as an exactly additive source}
Each row of the Qure.ai data holds a composite radiograph $T$ and the translated bone, soft-tissue and lung projections $B_t$, $S_t$ and $L$
(8-bit, $512\times512$, one sample per CT volume at 100~kVp). We verified that $T=a\,S_t+b\,B_t+c\,L+d$ holds up to rounding, with a per-sample bone
weight $b$ (\path{tools/qp_prep.py}); the additive bone image of the normalised input $X=T/255$ is therefore $B=b\,B_t/255$ and the remainder is
$S=X-B$. Over all \num{21887} rows the least-squares fit leaves a mean absolute residual of 0.36 grey levels per projection (largest 3.2 of 255).
The rows are split by patient into \num{21051} training, 406 validation and 430 test projections.

Training samples are composed as in Section~\ref{sec:methods}: bone re-mixing ($X'=S+wB$, exact), one random affine transform shared by all maps,
lesion insertion into $S$ with the bone target untouched, a monotone gamma tone curve ($X'=g(S+B)$, $S'=g(S)$, $B'=X'-S'$), linear blur applied to
$X'$ and $B'$, and noise and scatter applied to $X'$ only. Defining $B'=X'-S'$ keeps the decomposition exact under any monotone tone curve.

\subsection{Twin lesion-invariance regulariser}
Let $f$ be the suppressor, $x$ a training image and $b$ its exact bone image. A \emph{twin} of $x$ is $x^{+}=\min(x+n,1)$, where $n\ge0$ is the
signal of a few lesion patches taken from the lesion library of \emph{training} CTs and pasted at random positions in the lung fields, in composite
units and after all augmentation. The bone target of $x^{+}$ is the same $b$. Let $M$ be the lesion support dilated by $3$~px. The loss is
\begin{equation}
\begin{gathered}
\mathcal{L}=\tfrac12\big[\mathcal{L}_{\mathrm{sup}}(f(x),b)+\mathcal{L}_{\mathrm{sup}}(f(x^{+}),b)\big]+w\,\mathcal{L}_{\mathrm{inv}},\\
\mathcal{L}_{\mathrm{inv}}=\frac{\sum_p M_p\,\big|f(x^{+})_p-f(x)_p\big|}{\max\big(1,\sum_p M_p\big)},
\end{gathered}
\label{eq:twin}
\end{equation}
where $\mathcal{L}_{\mathrm{sup}}$ is the suppression loss of Section~\ref{sec:methods} and $w$ is the twin weight. Half of every batch consists of twins, so
the cost per step does not change. The supervised term already asks for the unchanged bone image on $x^{+}$; $\mathcal{L}_{\mathrm{inv}}$ ties the
two predictions directly on the lesion support. Each twin carries between one and two pasted lesions (uniform), each with a random scale and
amplitude; the supervised term is averaged over all base images and twins of the batch, and $\mathcal{L}_{\mathrm{inv}}$ is normalised per image by
its support and then averaged over the twins.

The ablation (Section~\ref{sec:experiments}) changes one part at a time: the twin weight $w$; a stop-gradient on $f(x)$; a support dilated further with a
$31\times31$ window, or the whole image, instead of $M$; up to four instead of up to two pasted lesions; and lesion insertion into the base image
without the tie term (the counterfactual insertion of Section~\ref{sec:methods}). All of these arms are in the ablation of
Section~\ref{sec:res_twin} (development data); the twin weights form the confirmatory dose-response.

\subsection{Deep ensemble and uncertainty}
The reported model is a deep ensemble of $K=3$ suppressors trained with seeds $0$, $1$ and $2$ \citep{lakshminarayanan2017ensembles}; seed $3$ is a reserve.
The bone image is the ensemble mean and the soft image is $x$ minus that mean. The ensemble standard deviation is the per-pixel spread. A split-conformal
scale $q$ is the $\lceil (n+1)(1-\alpha)\rceil/n$ quantile of $|b-\bar f|/(\sigma+\varepsilon)$ pooled over calibration pixels of the Qure-projection
validation CTs, so that $[\bar f-q(\sigma+\varepsilon),\,\bar f+q(\sigma+\varepsilon)]$ has nominal coverage $1-\alpha=0.9$
\citep{lei2018conformal,angelopoulos2023conformal}. Calibration treats pixels as exchangeable, so the guarantee is marginal over pixels, not
over images. The sparsification curve is the mean absolute error of the pixels that remain after removing the fraction $f\in[0,0.95]$ with the
largest spread; AUSE is the area between this curve and the oracle curve that removes pixels by their error (0 is perfect).

\subsection{Training budget}
The suppressor keeps the nested decoder and CBAM attention of Section~\ref{sec:methods} \citep{zhou2018unetpp,woo2018cbam}. The dose-response and
ablation arms use the SwinV2-T encoder of that section \citep{liu2022swinv2}; ours uses a ConvNeXt-V2-T encoder \citep{woo2023convnextv2}
(FCMAE pre-training, ImageNet-22k then ImageNet-1k fine-tuning \citep{deng2009imagenet}; \texttt{timm} \citep{wightman2019timm}
\path{convnextv2_tiny.fcmae_ft_in22k_in1k}) with a learning
rate $0.25\times$ that of the decoder. Every v2 run is trained at $512$~px with AdamW \citep{loshchilov2019adamw} (learning rate $2\times10^{-4}$,
batch 8), the schedule and loss of Section~\ref{sec:methods}, and \num{300000} samples (\num{37500} steps); evaluation uses the EMA weights. The
encoder screen uses \num{100000} samples (Section~\ref{sec:experiments}). The seeds per arm are listed in the generated tables.


\section{Exploratory screening rounds}
The operating point of the regulariser was found in a sequence of screening rounds. All of them are exploratory: single runs or few seeds, the
SZCH development pairs and the audit on all 75 SZCH images (450 sites), never a locked set. The numbers below are copied from the decision log
of the project; they are development evidence and are not used for any confirmatory claim.

\paragraph{Round 1: training source.} With SwinV2-T at \num{100000} samples, switching the source from LIDC renders to the Qure.ai projections
lowered the real-image absorption from 0.300 to 0.216, at a cost of about 0.5\,dB of development bone-region PSNR$_{cal}$; adding lesion
insertion on top (0.218) or mixing the two sources (0.286, 0.269 with insertion) did not do better, and a twin term of weight 1 on a mixed
source moved absorption only slightly (0.260). The reference model absorbed 0.134 on the same sites.

\paragraph{Round 2: twin weight.} On the Qure source alone, absorption fell monotonically with the twin weight: 0.216, 0.186, 0.151, 0.103 and
0.020 at $w=0,1,4,16,64$. Development bone PSNR$_{cal}$ stayed flat up to $w=16$ (21.69--21.79\,dB) and fell to 21.27\,dB at $w=64$; the bone
error on the Qure-projection validation data rose with the weight (MAE 0.0044 to 0.0082 at $w=16$). Seed repeats put the seed spread at about
0.02 in absorption and 0.1\,dB in PSNR; three times the training length and the larger SwinV2-B encoder each moved absorption by about one
seed spread. The twin weight was the lever.

\paragraph{Rounds 3 and 4: seeds, budget and the operating point.} At $w=64$ two seeds disagreed strongly (absorption 0.020 and 0.061,
PSNR 21.27 and 20.80\,dB), so that weight was not usable; $w=16$ to 32 was stable. At \num{300000} samples four seeds of SwinV2-T with
$w=16$ gave absorption 0.059--0.071 (mean 0.067) at a development PSNR of 21.73\,dB, tighter than at \num{100000} samples, while $w=32$ cost
about 0.4\,dB for a further 0.009 of absorption. The operating point was fixed at $w=16$.

\paragraph{Wave 1: dose-response and ablation at the full budget.} The twin weights $w\in\{0,4,8,16,32,64\}$ and the ablation arms were
trained at \num{300000} samples with at least two seeds; their development results fed no decision about the final configuration and are
reported only in the ablation table (Table~\ref{tab:ablation_twin}); the dose-response itself is confirmatory and is evaluated on the locked
sets below.

\paragraph{Encoder screen and the final configuration.} ConvNeXt-V2-T, DINOv3 ViT-S/16 and RAD-DINO ViT-B/14 were screened at
\num{100000} samples (Table~\ref{tab:app_screen}; rule R1, Chapter~\ref{ch:benchmark}). ConvNeXt-V2-T met R1 after a rerun with a reduced
encoder learning rate, RAD-DINO missed it on PSNR by 0.05\,dB but was continued as exploratory, and both met the replacement rule R2 at
\num{300000} samples. The pre-declared tie-break (mean absorption within 0.005, then the higher bone PSNR) chose ConvNeXt-V2-T at two seeds
each; with the third seeds the absorption gap became 0.0055, just outside the band, and the choice was not reopened. The final model is the
ensemble of three ConvNeXt-V2-T suppressors with $w=16$; SwinV2-T and RAD-DINO are reported as secondary encoders.

\section{Dose-response and ablation}
\label{sec:res_twin}
On the locked VinDr-RibCXR audit the absorption of the SwinV2-T family falls from 0.183 at $w=0$ to 0.043 at $w=64$
(Table~\ref{tab:dose_response}, Figure~\ref{fig:dose_response}). The slope of $A$ on $\log_2(1+w)$ is $-0.0233$ $[-0.0249,-0.0217]$; the
upper bound is below zero, so H1$'$ passes. On VinDr-RibCXR the fall is steepest at small weights and continues, more slowly, up to $w=64$; on
the Qure-projection test split with exact truth the slope is $-0.0174$ $[-0.0190,-0.0160]$ and absorption hardly changes above $w=16$ (0.118, 0.122 and 0.112). On the
held-out SZCH pairs the bone-region PSNR$_{cal}$ of the family stays between 21.22 and 21.28\,dB up to $w=16$ and drops to 20.92 and 21.01\,dB
at $w=32$ and 64 (Figure~\ref{fig:pareto}). The confirmatory data thus reproduce what the development data had shown: the regulariser lowers
absorption in proportion to the logarithm of its weight, and its fidelity cost appears only above $w=16$.

The ablation (Table~\ref{tab:ablation_twin}, development data) shows which parts carry the effect. A stop-gradient on the branch without the
pasted lesion removes it entirely ($A$ 0.185 against 0.067 for the full regulariser and 0.190 without it): the network must be able to move both
predictions towards each other. A consistency term over the whole image also removes it (0.197), presumably because the per-pixel weight of
the lesion is diluted by the much larger background. Dilating the support by 31\,px keeps part of the effect (0.137); up to four lesions per
twin instead of up to two keep all of it (0.074). Inserting lesions without the twin term (CLIA) gives 0.191, the same as no twin at all,
which matches the v1 finding on real images: showing the model lesions is not the same as telling it to leave them alone.
\begin{table}[htbp]
\centering\small
\caption{Dose-response of the lesion absorption on the twin weight (SwinV2-T family): seeds per weight, mean $A$ on VinDr-RibCXR and on the
Qure-projection test split, and the pre-specified slope with its image-cluster bootstrap interval.}
\label{tab:dose_response}
\fittable{tables/dose_response.tex}
\end{table}
\begin{figure}[htbp]
\centering
\safefig{figures/dose_response.pdf}
\caption{Dose-response of the lesion absorption $A$ on VinDr-RibCXR on the twin weight $w$ (SwinV2-T), one marker per seed, with the
pre-specified slope on $\log_2(1+w)$.}
\label{fig:dose_response}
\end{figure}
\begin{table}[htbp]
\centering\small
\caption{Ablation of the twin regulariser at $w=16$ (SwinV2-T, 300k samples, mean over seeds): one factor changed per row. Development evidence:
audit on all 75 SZCH images, bone PSNR$_{cal}$ on the 24 development pairs.}
\label{tab:ablation_twin}
\fittable{tables/ablation_twin.tex}
\end{table}

\section{Encoder results on the locked sets}
The twin effect does not depend on the encoder (Table~\ref{tab:encoders}). On VinDr-RibCXR it lowers absorption by 0.127 for ConvNeXt-V2-T,
0.114 for SwinV2-T and 0.112 for RAD-DINO, with intervals well away from zero, and on JSRT by 0.073, 0.061 and 0.064. RAD-DINO reaches the same
absorption as the selected ConvNeXt-V2-T on VinDr-RibCXR (0.053 against 0.051); single-seed results of all three families are in
Table~\ref{tab:app_seeds}. These comparisons are exploratory: the plan's confirmatory contrasts concern the selected model only.
\begin{table}[htbp]
\centering\small
\caption{Twin regulariser per encoder on the locked audit sets (exploratory): absorption difference to the matched no-twin control and its
interaction with bone overlap.}
\label{tab:encoders}
\fittable{tables/encoders.tex}
\end{table}
